\documentclass{article}

\usepackage[T1]{fontenc}
\usepackage{lmodern}
\usepackage{amsmath}
\usepackage{amssymb}
\usepackage{amsthm}
\usepackage{mathtools}
\usepackage[hidelinks]{hyperref}

\newtheorem{proposition}{Proposition}
\newtheorem*{proposition*}{Proposition}
\newtheorem*{conjecture*}{Conjecture}

\DeclareMathOperator{\wt}{wt}
\DeclareMathOperator{\tr}{tr}

\title{Planar 722: research outside the review draft: Recurrences, weight bounds, and deferred quantum interpretation}
\author{}
\date{Research notes, 11 November 2026}

\begin{document}

\maketitle

\textbf{Status and scope.} All derived claims are PROPOSITION, with the named gap
independent review\slash formalization pending. Distance 20 is CONJECTURE. None is THEOREM
(refereed complete). These notes are not part of the frozen DRAFT.tex review packet and are
not eligible for a final PROOF until fully refereed and formalized. Proposition numbering
\ref{prop:recurrence} and \ref{prop:one-block} is retained from the original draft to
preserve references in research discussions.

\textbf{Conventions.} Use the explicit finite model of DRAFT.tex: \(F = \mathbb{F}_2\),
\(Q = \{0, \dots, 18\} \times \{0, \dots, 18\}\), \(V = F^{722}\), \(f = x + x^2 + y^2\),
\(g = 1 + x^2 y + x^2 y^2\). The two qubit blocks are A and B. The row spaces of the
truncated X and Z check matrices are \(S_X, S_Z\). Weight counts physical coefficients;
\(P\) restricts to \(Q\).

\section{A new physical-weight recurrence}\label{sec:recurrence}

\setcounter{proposition}{2}
\begin{proposition}[forced propagation; independent review\slash formalization pending]%
  \label{prop:recurrence}
  For an X operator \((u, v)\) in \(V\), write \(u = \sum_{i=0}^{18} u_i(y) x^i\) and
  \(v = \sum_{i=0}^{18} v_i(y) x^i\). It is a cycle, meaning orthogonal to every Z row, if
  and only if
  \begin{equation}
    u_i = t u_{i-2} + F_i \quad (2 \le i \le 18),
  \end{equation}
  where
  \begin{equation}
    t = y + y^2, \quad F_i = v_{i-1} + v_{i-2} + y^2 v_i.
  \end{equation}
  In particular, for \(r = 0, 1\) and \(1 \le m \le \lfloor (18 - r)/2 \rfloor\),
  \begin{equation}
    u_{r+2m} = t^m u_r + \sum_{\ell=1}^{m} t^{m-\ell} F_{r+2\ell}.
  \end{equation}
\end{proposition}

\begin{proof}
  The overlap with the Z row at \((i, j)\) is the coefficient of \(x^i y^j\) in
  \(gu + fv\). Its possible second exponents all lie between 0 and 20, exactly the
  retained Z anchor range. Thus vanishing of the retained syndrome is equivalent to
  vanishing of entire polynomial slices for \(i = 2, \dots, 18\). Expanding that
  coefficient gives the first identity. For the second, the case \(m = 1\) is the first
  identity; substituting the formula for \(m - 1\) gives the stated sum, proving it by
  induction.
\end{proof}

\begin{proposition}[one-block weight bound; independent review\slash formalization pending]%
  \label{prop:one-block}
  Every nonzero X cycle supported on A alone has physical weight at least 29. This is a
  bound on all such cycles, without a logical nontriviality assumption. The value 29 is
  attained by a cycle.
\end{proposition}

\begin{proof}
  With \(v = 0\), the two parity channels are
  \begin{equation}
    u_{2m} = t^m u_0 \quad (0 \le m \le 9), \quad
    u_{1+2m} = t^m u_1 \quad (0 \le m \le 8).
  \end{equation}
  If \(u_0\) is nonzero, the degree of \(u_{18} = t^9 u_0\) is \(18 + \deg u_0\), by the
  elementary polynomial fact in DRAFT.tex. It is at most 18, so \(u_0 = 1\). If \(u_1\)
  is nonzero, the same reasoning at row 17 gives \(\deg u_1 \le 2\). Conversely these seed
  constraints suffice: all propagated slices have nonnegative degrees at most 18, and the
  displayed recurrence ensures the cycle equations.

  Multiplication by \(y^m\) does not change coefficient count, so it suffices to count
  \((1 + y)^m p\). Up to a monomial factor, the nonzero polynomials of degree at most two
  have exactly the four forms \(1\), \(1 + y\), \(1 + y^2 = (1 + y)^2\), and
  \(1 + y + y^2\). Indeed a one-term seed is a monomial; a two-term seed has exponent
  difference 1 or 2; a three-term seed must have all degrees 0,1,2. The relevant
  coefficient counts are as follows.

  \begin{center}
    \begin{tabular}{lllll}
      \(m\) & \(p = 1\) & \(p = 1 + y\) & \(p = 1 + y^2\) & \(p = 1 + y + y^2\) \\
      \hline
      0 & 1 & 2 & 2 & 3 \\
      1 & 2 & 2 & 4 & 2 \\
      2 & 2 & 4 & 2 & 4 \\
      3 & 4 & 2 & 4 & 4 \\
      4 & 2 & 4 & 4 & 6 \\
      5 & 4 & 4 & 8 & 4 \\
      6 & 4 & 8 & 2 & 6 \\
      7 & 8 & 2 & 4 & 8 \\
      8 & 2 & 4 & 4 & 6 \\
      Sum & 29 & 32 & 34 & 43 \\
    \end{tabular}
  \end{center}

  For completeness the table is a small polynomial identity calculation: in characteristic
  two \((1 + y)^{2^a} = 1 + y^{2^a}\), as follows by repeated squaring and cancellation of
  cross terms. Expressing \(m\) in binary therefore gives
  \((1 + y)^m = \prod_{a : \text{bit}_a(m) = 1} (1 + y^{2^a})\). Distinct subsets give
  distinct exponents, namely the integers whose binary bits are subsets of those of \(m\).
  This proves the first column. The next two columns use exponents \(m + 1\) and
  \(m + 2\). For the last column take the symmetric difference of that exponent set with
  its translates by 1 and 2. Explicitly its supports, for \(m = 0, \dots, 8\), are
\begin{verbatim}
{0,1,2}; {0,3}; {0,1,3,4}; {0,2,3,5};
{0,1,2,4,5,6}; {0,3,4,7}; {0,1,3,5,7,8};
{0,2,3,4,5,6,7,9}; {0,1,2,8,9,10}.
\end{verbatim}
  These explicit expansions verify every entry without any distance search. For the even
  channel one more slice is needed: at \(m = 9\) the seed 1 gives four coefficients. Thus
  its nonzero weight is \(29 + 4 = 33\). The two channels occupy disjoint rows, so their
  weights add. A nonzero odd channel has weight at least 29, and a nonzero even channel
  has weight 33. This proves the bound. Taking \(u_0 = 0, u_1 = 1\) attains 29, explicitly
  \begin{equation}
    u = \sum_{m=0}^{8} x^{1+2m} y^m (1 + y)^m, \quad v = 0.
  \end{equation}
  No claim that this cycle is a nontrivial logical is needed or made.
\end{proof}

\section{Distance-20 mission: what this checkpoint changes}\label{sec:mission}

\begin{conjecture*}
  Both nontrivial logical sectors of the explicit finite model have minimum weight at
  least 20. The main proposition in DRAFT.tex does not imply this.
\end{conjecture*}

Proposition~\ref{prop:one-block} excludes all A-only X cycles of weights 1 through 28.
Hence any hypothetical X logical of weight below 20 must involve the B block.
Proposition~\ref{prop:recurrence} identifies exactly how that block can suppress the
unforced growth: each \(v_i\) contributes \(y^2 v_i\) to \(F_i\) and contributes \(v_i\)
to each of \(F_{i+1}\) and \(F_{i+2}\) when those indices are retained. Cancellation must
be charged to these physical B coefficients; merely counting nonzero slices of the
homogeneous recurrence would be invalid.

The next mathematical target is a cancellation inequality for the forced recurrence,
modulo the X check space, with cost \(\sum_{i=0}^{18} (\wt(u_i) + \wt(v_i))\). It must
rule out nontrivial mixed-block cycles of cost at most 19, and an opposite-sector argument
is also needed. This target is open. It is more specific than a boundary-quotient
dimension calculation: the transfer multiplier, seed degree constraints, forcing
positions, and an exact unforced weight bound are now explicit. No exhaustive-search
certificate, UFD bulk argument, or Lean computation is part of this draft.

\subsection{Separation from input recognition}\label{sec:separation}

The definition and all propositions above stand independently of external data. The
companion geometry notes record a direct sparse match to the submitted 722 matrices and
list the upstream revision and input hash. That match is graded PROPOSITION with
independent replay pending. It is not a hypothesis of the main proposition in DRAFT.tex
and is not used in its proof. Likewise the submitted weight-20 witness is an upper bound
and supplies no part of the conjectured lower bound. This draft is an analytic checkpoint
toward distance 20, not a completed distance proof.

\section{Deferred quantum interpretation}\label{sec:quantum}

\begin{proposition*}[independent review and formalization bridge pending]
  Assuming the finite-model commutation and ranks from DRAFT.tex, the common positive
  stabilizer eigenspace has complex dimension \(2^8\). The proof below is preserved for
  later use. No formal bridge from the binary quotient to this Hilbert-space construction
  has been established in this task; it is therefore excluded from the review draft.
\end{proposition*}

For the quantum interpretation, take the complex Hilbert space with orthonormal basis
\(e_q\), indexed by \(q\) in \(F^{722}\). Define
\begin{equation}
  X(u) e_q = e_{q+u}, \quad Z(v) e_q = (-1)^{v \cdot q} e_q.
\end{equation}
These formulas imply \(Z(v) X(u) = (-1)^{u \cdot v} X(u) Z(v)\). Since the two row spaces
are orthogonal, the operators \(X(u) Z(v)\) for \((u, v)\) in \(S_X \times S_Z\) form a
commuting group under addition of labels. Their average
\begin{equation}
  E = 2^{-714} \sum_{u \in S_X} \sum_{v \in S_Z} X(u) Z(v)
\end{equation}
is a projection onto the common positive eigenspace: each group element fixes the
average, while on a vector fixed by all elements the average acts as identity. In
particular \(E^2 = E\). A nonidentity term has trace zero. If \(u \ne 0\), it permutes
every basis vector to a different one. If \(u = 0\) and \(v \ne 0\), pair basis labels
\(q\) and \(q + e_j\) at a coordinate where \(v_j = 1\); their diagonal signs cancel. The
identity term has trace \(2^{722}\). Thus \(\tr E = 2^8\). Finally any idempotent splits
the space as its image plus kernel, since \(w = Ew + (w - Ew)\) and their intersection is
zero; in the resulting basis its trace equals its image dimension. Trace is independent of
basis: direct summation gives \(\tr(AB) = \tr(BA)\), hence conjugation by any invertible
change-of-basis matrix preserves it. This proves the deferred quantum dimension
assertion. \qed

\end{document}
